\documentclass[10pt]{article}
\usepackage[margin=0.84in]{geometry}
\usepackage{amsmath,amssymb,mathtools}
\usepackage{booktabs,tabularx}
\usepackage[dvipsnames]{xcolor}
\usepackage[colorlinks=true,linkcolor=NavyBlue,urlcolor=BrickRed,hypertexnames=false]{hyperref}
\usepackage{microtype,fancyhdr,enumitem}
\definecolor{ink}{HTML}{19332D}
\definecolor{accent}{HTML}{187C68}
\definecolor{pale}{HTML}{ECF8F4}
\definecolor{passgreen}{HTML}{16794A}
\color{ink}
\pagestyle{fancy}\fancyhf{}
\lhead{Cyclic matrix groups}\rhead{Nonuniform $\Lambda$-distributions}\cfoot{\thepage}
\setlength{\headheight}{14pt}
\newcommand{\Sfun}{\mathcal S}
\newcommand{\Sym}{\operatorname{Sym}}
\newcommand{\Tr}{\operatorname{Tr}}
\newcommand{\PASS}{\textcolor{passgreen}{\textbf{PASS}}}
\newcommand{\summarybox}[1]{\begin{center}\fcolorbox{accent}{pale}{\parbox{0.91\linewidth}{#1}}\end{center}}
\title{\textbf{Nonuniform $\Lambda$-Distributions on Cyclic Matrix Groups}\\[3pt]
\large Fourier weights, convolution walks, and direct determinant checks}
\author{Proof and computational verification note}
\date{17 July 2026}

\begin{document}
\maketitle
\summarybox{\textbf{Outcome.}  The averaging-operator and convolution formulas
were checked on explicit diagonal models of $C_3,C_4,C_5$ in dimensions
$1,2,3$.  Arbitrary nonuniform measures and four walk times were tested for
every partition through total degree five.  All \textbf{288} comparisons pass;
the maximum absolute error is $2.91\times10^{-14}$.  Three additional
unexpanded determinant evaluations also pass.}

\section{The general coefficient theorem}
Let $G$ be finite, $\rho:G\to U(V)$, and let $\nu$ be a probability measure.
For finitely many active variables define
\begin{equation}
\Sfun_{\nu,V}(\mathbf t)=\sum_{g\in G}\nu(g)
 \prod_a\det(I-t_a\rho(g))^{-1}.                         \tag{1}
\end{equation}
The standard identity
$\det(I-tA)^{-1}=\sum_{r\geq0}\chi_{\Sym^r V}(A)t^r$
implies, by multiplying the series variable by variable, that
\begin{equation}
[\mathbf t^\tau]\Sfun_{\nu,V}
=\sum_g\nu(g)\chi_{W_\tau}(g)
=\Tr\!\left(\sum_g\nu(g)\rho_{W_\tau}(g)\right),
\qquad W_\tau=\bigotimes_a\Sym^{\tau_a}V.                \tag{2}
\end{equation}
This proves the proposed replacement of the Haar invariant dimension by the
trace of a nonuniform averaging operator.  Positivity of the measure does not
force the trace to be real or integral; the $C_5$ test below deliberately
exhibits a complex coefficient.

If $\nu=\kappa^{*r}$ and the convention is that a step product is
$X_1\cdots X_r$, then representation multiplicativity gives
\begin{equation}
\widehat\nu(W)=\widehat\kappa(W)^r.                        \tag{3}
\end{equation}
Equations (2)--(3) prove the convolution formula without a centrality
assumption.

\section{Specialization to \texorpdfstring{$C_m$}{C-m}}
Write $C_m=\{0,\ldots,m-1\}$ additively and
$\zeta_m=e^{2\pi i/m}$.  A diagonal representation is determined by weights
$w_1,\ldots,w_d\pmod m$:
\begin{equation}
\rho(k)=\operatorname{diag}(\zeta_m^{kw_1},\ldots,
\zeta_m^{kw_d}).                                           \tag{4}
\end{equation}
Let $M_\tau(q)$ denote the multiplicity of the character $k\mapsto
\zeta_m^{kq}$ in $W_\tau$.  Extracting weights from
\begin{equation}
\prod_{j=1}^d(1-xz^{w_j})^{-1}
=\sum_{r\geq0}\sum_q M_{(r)}(q)x^rz^q                    \tag{5}
\end{equation}
and convolving the multiplicities for the factors in $W_\tau$ gives the
finite Fourier formula
\begin{equation}
[\mathbf t^\tau]\Sfun_{\nu,V}
=\sum_{q\bmod m}M_\tau(q)\widehat\nu(q),
\qquad
\widehat\nu(q)=\sum_{k=0}^{m-1}\nu(k)\zeta_m^{kq}.        \tag{6}
\end{equation}
For a walk this becomes
\begin{equation}
[\mathbf t^\tau]\Sfun_{\kappa^{*r},V}
=\sum_qM_\tau(q)\widehat\kappa(q)^r.                     \tag{7}
\end{equation}
Thus the abstract matrix Fourier transform reduces to a transparent scalar
filter on each weight.

\section{Independent verification routes}
The direct route builds every matrix in (4), computes the power traces
$p_j=\Tr(\rho(k)^j)$, and obtains complete characters from Newton's recurrence
\begin{equation}
h_0=1,\qquad h_r=\frac1r\sum_{j=1}^r p_jh_{r-j}.            \tag{8}
\end{equation}
It then averages $\prod_a h_{\tau_a}$ using the fully enumerated
probability distribution.  The formula route never uses matrix traces: it
constructs the integer weight multiplicities in (5) and applies (6) or (7).
For walks, the direct distribution is obtained by repeated cyclic convolution,
whereas the formula route raises Fourier scalars to the $r$th power.

As a full target check, the code also evaluates (1) at
$(t_1,t_2)=(0.07,-0.04)$ using matrix determinants and compares it with the
product over the weights in (4).  These values stay safely inside every pole.

\section{Cases and results}
The representations and the nonuniform arbitrary measure are
\begin{center}
\begin{tabular}{@{}cccc@{}}
\toprule
Group & weights & dimension & $\nu(k)$ before normalization\\
\midrule
$C_3$ & $(1)$ & 1 & $1,2,3$\\
$C_4$ & $(1,-1)$ & 2 & $1,2,3,4$\\
$C_5$ & $(0,1,2)$ & 3 & $1,2,3,4,5$\\
\bottomrule
\end{tabular}
\end{center}
The walk step is $0.35\delta_0+0.40\delta_1+0.25\delta_{-1}$ and is tested
at $r=0,1,2,5$.  There are 19 partitions through total degree five, hence
$3\cdot(1+4)\cdot19=285$ coefficient comparisons plus three determinant
comparisons.

\begin{center}
\begin{tabularx}{\linewidth}{@{}lXrrc@{}}
\toprule
Case & coefficient & direct value & formula value & result\\
\midrule
$C_3$, arbitrary & $\tau=(3)$ & $1.0000000000$ & $1.0000000000$ & \PASS\\
$C_4$, walk $r=5$ & $\tau=(2,1)$ & $-0.0210525000$ & $-0.0210525000$ & \PASS\\
$C_5$, arbitrary & $\tau=(2,2)$ & $2.1666667+0.2293970i$ & $2.1666667+0.2293970i$ & \PASS\\
\bottomrule
\end{tabularx}
\end{center}
The complex $C_5$ row is a useful diagnostic: the implementation is genuinely
testing non-Haar averaging, not silently rounding to invariant dimensions.

\section{Reproducibility and scope}
Run \path{cyclic_groups/verification.py}, or open
\path{marimo_notebooks/nonuniform_cyclic.py} in marimo.  The notebook exposes the
degree cutoff and every comparison.  The finite computations verify all
coefficients in the stated range; equations (1)--(7) are all-degree proofs.
No sampling is used.
\end{document}
